\section{Evaluation}
\label{sec:evaluation}

We organize the evaluation around four questions: (1) can software change the
DUT-visible subsystem without changing its bitstream, (2) do native drivers
remain functional, (3) what control and data overhead does mediation introduce,
and (4) how does provisioned endpoint capacity affect FPGA cost?

\subsection{Experimental Setup}

The DUT is a XiangShan RISC-V SoC prototype running Linux at 50 MHz. The host
communicates with the front-end through XDMA and runs the \sys QEMU back-end.
The evaluated devices include physical NVMe SSDs and an Intel 82580 Ethernet
function. NVMe uses Linux \texttt{fio}; network validation uses DHCP, ping, and
\texttt{iperf}. \TODO{Add FPGA board/part, host CPU, Linux/QEMU versions, XDMA
version, exact device models, link widths/speeds, fio parameters, iperf version,
run count, and error bars.}

\subsection{Software-Defined Composition and Compatibility}

With the same physical wiring and FPGA image, changing only the backend
configuration exposes different populated endpoint sets under the fixed switch
capacity. For each configuration, Linux discovers the bridges and endpoints,
assigns BAR addresses, binds its unmodified \texttt{nvme} or \texttt{igb}
driver, and issues I/O through standard applications. The NVMe path completes
queue initialization and \texttt{fio} traffic. The NIC path reaches physical
link and carrier, completes DHCP in both directions, installs the lease, and
passes bidirectional ping. These results establish a stronger property than
enumeration: configuration, control, translated DMA, completion, and interrupt
handling form a working loop.

\TODO{Add a table whose rows are tested software configurations and whose
columns report topology, backend location, driver bind, initialization, and
workload result. Include an NVMe+NIC simultaneous run to demonstrate composition
rather than isolated device support.}

\subsection{Control-Path Overhead}

Table~\ref{tab:control} separates a basic BAR write from an effective submission
that includes device-specific preparation before the physical doorbell. Host-
mediated control takes tens of microseconds. The submission path costs more than
a simple write because the adapter waits for stable queue memory, translates
DMA references, and preserves the submission ordering invariant.

\begin{table}[t]
\centering
\caption{Measured host-mediated control latency from the current experiment.}
\label{tab:control}
\begin{tabular}{@{}lrr@{}}
\toprule
Device & BAR write (\si{\micro\second}) & Effective doorbell (\si{\micro\second}) \\
\midrule
NVMe & 30.15 & 61.41 \\
Intel NIC & 6.20 & 44.42 \\
\bottomrule
\end{tabular}
\end{table}

\TODO{Define both latency endpoints precisely, reconcile these wall-clock
measurements with the cycle-count results in the presentation, add distributions
and repetitions, and decompose FPGA transport, QEMU dispatch, semantic
translation, and physical BAR access.}

\subsection{End-to-End NVMe Performance}

Table~\ref{tab:nvme} compares two physical placements already measured through
the DUT. Sequential throughput is close: the on-board path reaches 98.2\% of the
host-attached sequential-read result and slightly exceeds its sequential-write
result. Random-read IOPS are also close. Random-write results vary substantially
and require repeated trials before supporting a general claim.

\begin{table}[t]
\centering
\caption{Initial NVMe results observed by the DUT.}
\label{tab:nvme}
\begin{tabular}{@{}lrr@{}}
\toprule
Workload & Host-attached & On-board \\
\midrule
128-KiB sequential read (MiB/s) & 13.52 & 13.28 \\
128-KiB sequential write (MiB/s) & 12.34 & 12.91 \\
4-KiB random read (IOPS) & 176.41 & 174.50 \\
4-KiB random write (IOPS) & 140.62 & 178.07 \\
\bottomrule
\end{tabular}
\end{table}

At 50 MHz and queue depth one, end-to-end throughput is dominated by the slow
DUT and its software submission path. The current data show that tens of
microseconds of control mediation do not dominate these particular workloads;
they do not establish that the overhead will remain hidden for a faster DUT or
higher queue depth. \TODO{Add block-size and queue-depth sweeps, direct-device
baselines, CPU utilization, and repeated measurements with confidence
intervals.}

\subsection{Network Functionality and Performance}

The single-queue NIC backend achieves 56.65 Mbit/s for TCP from the DUT to a
peer and 46.82 Mbit/s in the reverse direction. Measured UDP rates are 4.90 and
7.75 Mbit/s, and ping RTT is 2.80 ms. These values demonstrate a functioning
translated TX/RX path, not line-rate Ethernet. \TODO{Explain the unusually low
UDP results, add loss rate and datagram size, run repeated single-stream TCP
tests, and compare with a direct or on-board NIC baseline.}

\subsection{FPGA Cost and Scaling}

The central hardware claim requires synthesis at multiple provisioned
capacities. \TODO{Report LUT, FF, BRAM/URAM, and Fmax for 1, 2, 4, 8, and 13
endpoint slots; separate fixed front-end cost, configuration-shadow capacity,
and per-route cost. Compare against the earlier single-device proxy or another
FPGA-heavy baseline. Also report RTL changes and resynthesis requirements when
adding the NIC adapter.}

The expected result is not zero per-endpoint cost. Rather, the FPGA slope should
reflect generic configuration and routing capacity, while device-specific queue
and protocol state resides mainly in software. We will state this conclusion
only after the synthesis data are included.
